\lecture{10}{Thu 22 Sep 2022 12:07}{}
\subsection{Cauchy's Theorem}

\begin{theorem}[Cauchy's Theorem]
	If $f$ is analytic in a region $\Omega$ and $\overline{\mathcal{D}} =  \mathcal{D} \cup \partial \mathcal{D}\subset \Omega$ , $\partial \mathcal{D}$ is piecewise smooth curve, then \[
		\int _{\partial \mathcal{D}} f(z) dz =0
	.\]  
\end{theorem}

\subsubsection{Continuous Deformation of curves}

\begin{definition}
	Consider a curve $\gamma_0(t): [a,b] \to \Omega$. Consider a family of curves, $\{\gamma_s(t): [a,b] \to  \Omega\} $, such that $a\le t\le b$ and $0\le s\le 1$, $\gamma_s(a) = \gamma_0(a)$, $\gamma_s(b) = \gamma_0(b)$.

	Then $\gamma_s$ is a \textit{continuous deformation} from $\gamma_0$ to $\gamma_1$.
\end{definition}

\begin{corollary}
	If $f$ is analytic in $\Omega$, and $\gamma_0$ and $\gamma_1$ are two curves in $\Omega$ which can be continuously deformed, then \[
		\int _{\gamma_0}f(z) dz
		=
		\int _{\gamma_1}f(z) dz
	.\] 
\end{corollary}
\begin{proof}
	Let $\mathcal{D}$ be the region bounded by $\gamma_0$ and $\gamma_1$. Now $\partial \mathcal{D} = \gamma_0 - \gamma_1$, and the integral in $\mathcal{D}$ is zero.
\end{proof}

\begin{definition}[Simply Connected]
	A region $\mathcal{D}$ is called \textit{simply connected} if any two curves with common begin and common end in this region can be deformed to each other.
\end{definition}

\begin{proposition}
	A region is simply connected iff 
	every closed curve can be deformed to a point.
\end{proposition}

\begin{intuition}
	Simply connected means the region has "no holes".
\end{intuition}

\begin{example}
	\begin{enumerate}
		\item Every conved region is simply connected.
		\item Every starlike region is simply connected.
		A "starlike" region is a region where there exists a point $z_0$ such that every light segment from this point to some other point in the region is contained in the region.
		Define \[
			\gamma_s(t) = z_0(1-s) + s\gamma(t)
		.\] 
	\item $\mathbb{C}\setminus  (-\infty , 0]$ is not convext, but starlike, $z_0 = 1$.
	\item A branch of $\log z$ exists in $\mathcal{D}$ if and only if $\mathcal{D} \subset \Omega$, where $\Omega$ is simply connected, $0 \not\in \Omega$.
	\end{enumerate}
\end{example}

\begin{theorem}(Cauthy's Theorem for simply connected regions)
	In a simply connecteed region, for any analytic funtion $f$, integrals \[
		\int _{\gamma}f(z) dz
	\] are path independent. 
\end{theorem}

\begin{example}[Winding Number]
	\begin{enumerate}
\item The integral $\int _{\gamma}\frac{dz }{z}$ is $2 \pi i$ times the \textit{winding number} of the curve around origin. The integral only depends on the times the cirve circulates around origin.
\item The integral \[
	\frac{1}{2 \pi i} \oint _{\gamma} \frac{dz}{z - a}
.\] 
is the \textit{winding number } about $a$.

Intuitively, it is the number the cirve circulates around $a$ \textbf{counter-clockwise}.
\item \[
	\oint_{\gamma}\frac{1}{\left( z-a \right) ^2}dz = 0
.\] Reason: the function has a primitive $- \frac{1}{z - a}$. So we can ignore the terms with order larger than $1$ in the partial fraction decomposition.
\end{enumerate}
\end{example}

\begin{theorem}
	In a simply connected region, every real harmonic function $u$ is the real part of some analytic function $f$ in the same region (also imaginary part).
\end{theorem}

\begin{enumerate}
	\item Complex conjugate of gradient of harmonic function is always analytic.
	\item In a simply connected region, $g$ has a primitive ($g$ is analytic and region is simply connected, therefore $g$ has a primitive), $f' = g$, $f$ is the analytic function we are looking for.
		\[
			f = p + i q, \quad f' = p_x + i q_x = p_x - i p_y = g = u_x - i u_y
		.\] 
		So $p_x = u_x, p_y = u_y$. So \[
			\operatorname{Re} f = p = u + C
		.\] 
\end{enumerate}
